% Part: history
% Chapter: biographies
% Section: ernst-zermelo

\documentclass[../../../include/open-logic-section]{subfiles}

\begin{document}

\olfileid[pt]{his}{bio}{zer}

\olsection{Ernst Zermelo}

\olphoto{zermelo-ernst}{Ernst Zermelo}

Ernst Zermelo nasceu em 27 de julho de 1871 em Berlim, na Alemanha.
Tinha cinco irmãs, embora a família sofresse de frágil saúde e apenas
três sobrevivessem à idade adulta. Os pais também faleceram quando ele
era jovem, deixando-o órfão, com os irmãos, aos dezessete anos.
Zermelo tinha um profundo interesse pelas artes, e sobretudo pela
poesia. Era conhecido por ser perspicaz, espirituoso e crítico. Entre
suas conquistas matemáticas mais celebradas estão a introdução do
axioma da escolha (em 1904) e sua axiomatização da teoria dos
conjuntos (em 1908).

Os interesses de Zermelo na universidade eram variados. Frequentou
cursos de física, matemática e filosofia. Sob orientação de Hermann
Schwarz, Zermelo concluiu a dissertação \emph{Investigations in the
  Calculus of Variations} em 1894 na Universidade de Berlim. Em 1897,
decidiu prosseguir estudos na Universidade de G\"{o}ttingen, onde foi
fortemente influenciado pelo trabalho fundacional de David Hilbert.
Em 1899 tornou-se elegível para uma cátedra, mas só obteve uma onze
anos depois---possivelmente devido a seu comportamento estranho e à
``pressa nervosa.''

Zermelo obteve finalmente uma cátedra remunerada na Universidade de
Zurique em 1910, mas foi forçado a se aposentar em 1916 por causa da
tuberculose. Depois da recuperação, foi nomeado professor honorário na Universidade
de Freiburg em 1921. Nesse período
dedicou-se à matemática fundacional.  Irritou-se com os trabalhos de
Thoralf Skolem e Kurt G\"{o}del, e criticou publicamente suas
abordagens em seus artigos.  Foi demitido do cargo em Freiburg em
1935, devido a sua impopularidade e à oposição à ascensão de Hitler
ao poder na Alemanha.
 
Os últimos anos da vida de Zermelo foram marcados pelo isolamento.
Depois da demissão em 1935, abandonou a matemática. Mudou-se para
o campo, onde viveu modestamente. Casou-se em 1944 e ficou completamente
dependente da esposa à medida que ia ficando cego. Zermelo perdeu
totalmente a visão por volta de 1951. Faleceu em G\"{u}nterstal,
na Alemanha, em 21 de maio de 1953.

\begin{reading}
Para uma biografia completa de Zermelo, ver \citet{Ebbinghaus2015}.
Os artigos fundadores de 1904 e 1908 de Zermelo podem ser lidos no
alemão original \citep{Zermelo1904,Zermelo1908}.  As obras reunidas
de Zermelo, incluindo os escritos sobre física, estão disponíveis em
tradução inglesa em \citet{Ebbinghaus2010,Ebbinghaus2013}.
\end{reading}

\end{document}
